\begin{exercise}[subtitle={Stochastic Simulation and Quantiles \Coffeecup\Coffeecup}]
    \label{ex: stochastic simulation}
    The goal is to simulate a random variable \(X\) with distribution 
    function \(F(x) = \prob{X \le x}\) using a standard uniform random variable
    \(U\sim \uniform(0,1)\). We define the generalized inverse of \(F\) as
    \[
        F^{\leftarrow}(p) \coloneq \inf\{x \in \real : F(x) \ge p\}, \quad p \in (0,1).
    \]
    Show that
    \begin{enumerate}[label=(\alph*), noitemsep]
        \item\label{it: quantile if invertible}
        if \(F\) is invertible, then \(F^{\leftarrow}(p) = F^{-1}(p)\) for all
        \(p \in (0,1)\),

        \item\label{it: monotonicity}
        \(F^{\leftarrow}\) is non-decreasing,

        \item\label{it: F-(F(x)) le x}
        \(F^{\leftarrow}(F(x)) \le x\) for all \(x \in \real\),

        \item\label{it: F(F-(p)) > p}
        \(F(F^{\leftarrow}(p)) \ge p\) for all \(p \in (0,1)\),

        \item\label{it: F-(p)< x iff p<F(x)} \(F^{\leftarrow}(p) \le x\) if and only if \(p \le F(x)\),

        \item\label{it: inverse transform sampling}
        \(X\coloneq F^{\leftarrow}(U) \sim F\) for a standard uniform \(U\sim \uniform(0,1)\), i.e.\ 
        \[
            F(x) = \prob{X \le x}
            \qquad \forall x \in \real.
        \]

        \item\label{it: simulate exponential rvs} Propose a way to simulate an exponential distribution \(X \sim
        \exponential(\lambda)\) using a standard uniform random variable \(U\sim
        \uniform(0,1)\).
    \end{enumerate}
    The function \(F^{\leftarrow}\) is also called a \emph{quantile function} of
    the distribution function \(F\). The reason for this is
    the following: Show that
    \begin{enumerate}[resume, label=(\alph*), noitemsep]
        \item\label{it: quantiles}
         \(Q(p) \coloneq (-\infty, F^{\leftarrow}(p)]\) is the smallest
        interval of the form \((-\infty, a]\) such that \(\prob{X \in (-\infty,
        a]} \ge p\) for \(X\sim F\). We say that \(Q(p)\) is a \(p\)-quantile.
    \end{enumerate}
\end{exercise}
\begin{solution}
\begin{enumerate}[wide=0pt]
    \item[\ref{it: quantile if invertible}]
    Since \(F\) is monotonically increasing as a distribution
    function and additionally invertible, it is strict monotonically
    increasing. That is
    \[
        F(x) < F(y) \quad \text{for all } x < y.
    \]
    So clearly for all \(x< F^{-1}(p)\) we have \(F(x) < p\) and
    therefore they are not contained in the set \(\set{x\in \real : F(x)
    \ge p}\). This implies \(F^{-1}(p)\) is the smallest element in this set
    and consequently \(F^{\leftarrow}(p) = F^{-1}(p)\).

    \item[\ref{it: monotonicity}]
    Let \(p < q\). Then for any \(x \in \real\) such that \(F(x) \ge q\)
    we also have \(F(x) \ge p\). Consequently,
    \[
        \set{x \in \real : F(x) \ge q} \subseteq \set{x \in \real : F(x) \ge p},
    \]
    This implies the infimum of the former is greater than the infimum of the latter
    and thus
    \[
        F^{\leftarrow}(p) = \inf\{x \in \real : F(x) \ge p\} \le \inf\{x \in \real : F(x) \ge q\} = F^{\leftarrow}(q).
    \]

    \item[\ref{it: F-(F(x)) le x}]
    Since \(x\in \set{y\in \real : F(y) \ge F(x)}\), we have
    \[
        x \ge \inf\set{y\in \real : F(y) \ge F(x)} = F^{\leftarrow}(F(x)).
    \]

    \item[\ref{it: F(F-(p)) > p}]
    Let \(x_n \in \set{x\in \real : F(x) \ge p}\) be a sequence with \(x_n \downarrow F^{\leftarrow}(p)\). Then
    by right-continuity of the distribution function \(F\)
    \[
        p \le \lim_{n\to \infty} F(x_n)
        = F\Bigl(\lim_{n\to \infty} x_n\Bigr)
        = F(F^{\leftarrow}(p)).
    \]
    
    \item[\ref{it: F-(p)< x iff p<F(x)}]
    Assume \(F^{\leftarrow}(p) \le x\). Then by \ref{it: F(F-(p)) > p} 
    and the monotonicity of \(F\) we have
    \[
        p \overset{\ref{it: F(F-(p)) > p}}\le F(F^{\leftarrow}(p)) \le F(x).
    \]
    For the opposite direction assume \(p\le F(x)\), then by the
    monotonicity of \(F^{\leftarrow}\) \ref{it: monotonicity}
    and \ref{it: F-(F(x)) le x} we have
    \[
        F^{\leftarrow}(p) \overset{\ref{it: monotonicity}}\le F^{\leftarrow}(F(x)) \overset{\ref{it: F-(F(x)) le x}}\le x.
    \]

    \item[\ref{it: inverse transform sampling}]
    We simply use \ref{it: F-(p)< x iff p<F(x)} and the definition of \(U\) to get
    \[
        \prob{X \le x}
        = \prob{F^{\leftarrow}(U) \le x}
        \overset{\ref{it: F-(p)< x iff p<F(x)}}\le \prob{U \le F(x)}
        = \int_0^{F(x)} 1 \, du
        = F(x).
    \]

    \item[\ref{it: simulate exponential rvs}]
    The CDF of an exponential distribution is \(F(x) = 1 - e^{-\lambda x}\) for \(x \ge 0\).
    Since this function is invertible on \([0, \infty]\) we use
    \ref{it: quantile if invertible}. To get
    \[
        F^{\leftarrow}(p) = -\tfrac{1}{\lambda} \log(1-p).
    \]
    Using \(U \sim \uniform(0,1)\) we thus set \(X = F^{\leftarrow}(U)\),
    and finish by \ref{it: inverse transform sampling}.

    \item[\ref{it: quantiles}]
    For any \(a < F^{\leftarrow}(p)\) we have by definition of \(F^{\leftarrow}\) that
    \(F(a) < p\) and thus
    \[
        \prob{X \in (-\infty, a]} = F(a) < p.
    \]
    However for \(a=F^{\leftarrow}(p)\) we have \(F(a)\ge p\) by
    \ref{it: F(F-(p)) > p}.  Consequently, \(F^{\leftarrow}(p)\) is the
    smallest \(a\) such that \(\prob{X \in (-\infty, a]} \ge p\).
    \qedhere
\end{enumerate}
\end{solution}
